\chapter*{Abstract}
\addcontentsline{toc}{chapter}{Abstract}
This report details the architectural analysis and design of the Billing Service within the ComOps Enterprise Resource Planning ecosystem. The system originally maintained local duplicates of core database tables, causing data inconsistencies and replication of business logic. To resolve these integrity issues, we present a refactoring plan that migrates the service to a fully reactive Hexagonal Architecture using Spring Boot WebFlux and WebClient. This allows the Billing Service to dynamically fetch client and supplier data from the central Kernel module, which acts as the single source of truth. The multi-tenant context is propagated reactively through HTTP headers and Reactor context. Our results demonstrate that this integration eliminates data duplication, ensures transactional integrity, and preserves reactive non-blocking execution performance. Future steps include testing network resilience under load and implementing circuit breakers.

\chapter*{Résumé}
\addcontentsline{toc}{chapter}{Résumé}
Le présent rapport expose le travail d'analyse et de conception technique réalisé pour optimiser l'intégration du module de facturation au sein du progiciel de gestion intégré ComOps. Dans sa version initiale, le service de facturation souffrait d'une duplication locale des données relatives aux clients et aux fournisseurs, ce qui générait des incohérences de données régulières et nuisait à l'intégrité du système d'information. Afin de remédier à ces dysfonctionnements, nous avons conçu et mis en œuvre une refonte architecturale s'appuyant sur les principes de l'architecture hexagonale et le modèle de programmation réactif de Spring Boot. Cette approche permet de déléguer la gestion des référentiels au noyau central du système, en remplaçant la persistance locale par des requêtes réseau non bloquantes et réactives pour interroger dynamiquement les données, tout en assurant l'isolation des organisations par la propagation automatique de l'identifiant d'organisation dans le contexte d'exécution. Les résultats de cette refonte montrent une centralisation complète des données de tiers, une suppression des tables locales redondantes et le maintien de la non-régression sur l'ensemble des processus de facturation et de génération de documents. Bien que cette solution améliore la cohérence globale, elle introduit une dépendance forte vis-à-vis de la disponibilité réseau du noyau central, ce qui ouvre des perspectives de recherche sur l'intégration de disjoncteurs pour sécuriser le fonctionnement en cas de panne réseau temporaire.

\vspace{1cm}
\noindent \textbf{Mots-clés :} Facturation, Noyau Central, Architecture Hexagonale, Programmation Réactive, Multi-tenant, Intégration de Services, Base de Données Relationnelle Réactive.
